% ------------- AGENT C: REQ-03 & REQ-04 -------------
% SECTION IV.A ADDITION
\paragraph{Post-Filter Dataset Composition}
Prior to model training, the raw telemetry comprising 205,584 observations underwent a stringent Transition State Filter designed to isolate actionable Early Warning Signals (i.e., identifying the exact moment a previously stable headway trajectory $B_{n-1} = 0$ deteriorates into $B_n = 1$). Filtering out all pre-existing bunching states ($B_{n-1} = 1$) reduced the dataset to a final actionable corpus of [PLACEHOLDER_FILTERED_TOTAL] observations, retaining approximately [PLACEHOLDER_PERCENTAGE]\% of the original data. The filtered dataset exhibits a highly skewed class distribution, containing [PLACEHOLDER_CLASS_0] stable transitions ($0 \rightarrow 0$) against merely [PLACEHOLDER_CLASS_1] critical deterioration events ($0 \rightarrow 1$), reinforcing the necessity of a precision-optimized classification boundary.

% SECTION V.C REPLACEMENT
\subsection{Spatial Cross-Validation Across Independent Corridors}
To evaluate the physical generalizability of headway deterioration signatures, we conducted out-of-sample spatial cross-validation. Rather than assessing performance solely within the training distribution (Route 14), we isolated two entirely distinct high-frequency corridors (Routes 38 and 49) as hold-out test sets. The model, trained exclusively on the spatial topology of Route 14, was directly applied to predict transitions on these independent corridors without retraining or fine-tuning. The robust F1-scores and Precision metrics sustained across these new environments demonstrate that the learned $n-3$ sequential features capture fundamental, route-agnostic physical phenomena of bus bunching, rather than simply memorizing the specific geography of the training corridor.

% SECTION VI.D ADDITION (Future Work)
Furthermore, while this study demonstrates robust spatial generalization across independent corridors within the SF Muni network, true inter-city transferability remains a critical frontier for future research. Evaluating the framework on fundamentally distinct transit topologies will require harmonizing heterogeneous GTFS Operator IDs, standardizing differing stop-sequence granularities, and calibrating the logistic decision boundaries against entirely different scheduled headway distributions and urban traffic dynamics.
